\documentclass[11pt]{article}
\usepackage[letterpaper,margin=1in]{geometry}
\usepackage[T1]{fontenc}
\usepackage{lmodern,amsmath,amssymb,booktabs,microtype}
\usepackage[hidelinks]{hyperref}
\newenvironment{proof}{\par\noindent\textit{Proof.} }{\unskip\nobreak\hfill\hbox{$\square$}\par}
\newtheorem{theorem}{Theorem}
\newtheorem{lemma}[theorem]{Lemma}
\newcommand{\result}{11.6854827408}
\newcommand{\wnum}{4581442652581}
\newcommand{\wden}{392062763189}
\hypersetup{pdftitle={A certified lower bound for the growth constant of plane meanders},pdfauthor={Iaroslav Argunov}}
\title{\bfseries A certified lower bound for the growth constant of plane meanders}
\author{Iaroslav Argunov\\{\small Independent Researcher, Oakland, California, USA}}
\date{November 22, 2026}
\begin{document}
\maketitle
\begin{abstract}
Let \(M_n\) count the connected closed plane meanders on \(2n\) ordered sites. We prove that \(\lim M_n^{1/n}>\result\) by confining the usual connectivity transfer to at most 32 exposed ends. The certificate compares integer vectors coordinate by coordinate. A separate computation using only returns to one state gives \(11.421465\) as a strict lower bound at width 26. Both improve the published lower bound \(11.380\). A deterministic relative approximation with prescribed error \(\epsilon\) takes \(O^*(2^{4n/7}\epsilon^{-3/2})\) bit operations. The programs and integer comparisons accompany the paper.
\end{abstract}

\section{Counting convention and result}
Fix the order of sites on a horizontal line. A closed meander is an ordered pair of noncrossing perfect matchings on the sites, drawn on opposite banks, whose union is one cycle. In particular \(M_1=1,\ M_2=2,\ M_3=8\). Albert and Paterson proved the existence of
\(\mu=\lim_{n\to\infty}M_n^{1/n}\) and obtained \(11.380\le\mu\le12.901\) by bounding the language of meander words \cite{ap}. Caroppo, Da Lozzo, Di Battista and Vihrovs give the upper bound \(12.806\) and recall \(11.380\) on the lower side \cite{caroppo}. Borga, F\'eray and Winkler record the interval \([11.380,12.806]\) in their recent survey \cite[p.~41]{bfw}. The constants here use the order \(n\); with the number of crossings as exponent their square roots apply.

\begin{theorem}\label{main}
The ordinary meandric growth constant satisfies
\[
             \mu\ \ge\ \frac{\wnum}{\wden}\ >\ \result .
\]
\end{theorem}

We use Jensen's transfer through a vertical cut \cite{jensen}. The number of exposed ends is bounded at every cut, however long the final word. One finite transfer can therefore supply a bound for arbitrarily large orders. Restricting allowed configurations to obtain growth estimates has precedents in Bacher's discussion of exterior arches \cite{bacher} and in the sublanguages of Albert and Paterson. Our contribution is the explicit integer inequality for this strip, together with a direct return calculation that already crosses the earlier bound. The argument uses no estimated critical exponent.

\newpage
\section{A finite transfer inside the connected language}
Order the \(2k\) exposed ends by following the boundary of the scanned half-plane from the uppermost to the lowermost end. Write \((\pi,u)\) for their noncrossing pairing by paths already drawn and the number \(u\) of upper ends. The slit lies between ends \(u\) and \(u+1\), with \(0\le u\le 2k\). Retain \(1\le k\le W/2\) for an even \(W\ge2\). The first site necessarily opens two ends and reaches \(o=(\{1,2\},1)\). Strip off this site and the final site, which closes its two ends in state \(o\).

At an interior site there are at most four choices. Opening on both banks inserts an adjacent pair at the slit and increases \(u\) by one. Closing on both deletes the two ends next to the slit and joins their former partners, decreasing \(u\) by one. This move requires ends on both banks that belong to different paths. Closing one bank and opening the other moves the slit one place towards the bank being closed and preserves \(\pi\). Omit moves that leave the strip or pop an empty bank. Give each remaining choice weight one, including its multiplicity if two choices reach the same stored state. Let \(T_W\) be the resulting matrix, with source rows. Its two-step restriction \(A_W=T_W^2|_{u\ {\rm odd}}\) contains the root.

\begin{lemma}\label{strip}
Writing \(C_n(W)\) for the connected meanders of order \(n\) confined to this strip, one has \(C_n(W)=(A_W^{\,n-1})_{oo}\). Every physical state communicates with \(o\) within its own width. Moreover
\[
 \lambda_W:=\lim_{n\to\infty} C_n(W)^{1/n}=\rho(A_W),\qquad
        \lambda_W\uparrow\mu\quad (W\to\infty,\ W\ {\rm even}).
\]
\end{lemma}
\begin{proof}
The two bracket words determine the arch matchings uniquely. The closing rule excludes exactly the closures of a component before the final site. Appending the compulsory last site to a return closes the sole remaining path, proving the counting identity.

Shifts let us place the slit at any gap without changing the pairing. For \(2k\ge4\), some two consecutive ends are not partners (the two neighbors of an interior end cannot both be its partner). Shift to such a gap and apply the double close. Repetition reaches width two, where a shift reaches \(o\). In the reverse direction every noncrossing pairing has an adjacent paired pair: take a shortest arc. Shift the slit between its ends and undo a double opening. Repetition constructs a path from \(o\) to the original state. The widths never exceed the initial width in either reduction. Every move flips the slit parity; at \(o\) a shift immediately followed by its reverse is legal. Thus the two-step matrix on the odd class is primitive, giving the asserted growth of returns.

The counts increase with \(W\) and are bounded above by \(M_n\). At \(W=2n\) every order-\(n\) word fits, since at site \(t\) there are at most \(2\min(t,2n-t)\) exposed ends. Concatenating \(h\) interior returns of \(n-1\) slots yields \(M_n^h\) distinct returns in \(h(n-1)\) slots. Hence \(\lambda_{2n}\ge M_n^{1/(n-1)}\) for each \(n\ge2\). Taking limits proves the last assertion.
\end{proof}

\newpage
\section{Integer witnesses}
For a nonnegative matrix \(A\) and a vector \(v\) with all entries positive, put \(L(v)=\min_i(Av)_i/v_i\). Positivity implies \(A^m v\ge L(v)^m v\) entrywise. Thus \(L(v)\le\rho(A)\). One can avoid spectral theory altogether. If every state reaches \(o\) in exactly \(b\) two-step slots, then \(A^b e_o\ge{\bf1}\). Consequently
\[
 (A^{m+b})_{oo}\ \ge\ (A^m v)_o/\max_i v_i
             \ \ge\ L(v)^m\,v_o/\max_i v_i .
\]
The existence of a common \(b\) follows from Lemma~\ref{strip} by padding at \(o\). For the stored graphs the programs also verify such a \(b\) directly, starting with the root indicator. This argument only needs the count of root returns and a return from each stored state.

\paragraph{Representing the two banks.}
At width \(2k\), reflection swaps the banks, reverses the order of the ends, and replaces \(u\) by \(2k-u\). Store all pairs with \(u\le k\), keeping every pairing at \(u=k\). At this boundary the shift to \(k+1\) is evaluated at the reflected pairing and \(k-1\). A double opening or closing stays in the stored half after the change of width. Both shift choices are counted even when they have the same destination. If a full vector is invariant under bank exchange, one stored sweep is precisely the restriction of the full sweep. Write \(\widetilde A_W\) for the stored two-step matrix. Induction from the invariant root indicator gives \((\widetilde A_W^m e_o)_o=(A_W^m e_o)_o\) for every \(m\ge0\). Below, \(A\) denotes the matrix used for that comparison. The number of stored cells is \(\sum_{k=1}^{W/2}(k+1)\mathrm{Cat}_k\). In particular no orbit sizes divide the count.

\paragraph{Constructing the witness.}
Start with \(2^{52}\) in each odd cell. Form the unrounded \(z=Av\) by an even sweep and an odd sweep. At a recorded round find the smallest ratio \(z_i/v_i\) by integer cross multiplication over all odd cells, checking positivity. Between rounds rescale by a common positive factor and round down; the choice serves only to find a vector with a good minimum. The row has at most four letters per site, so \(\max z\le16\cdot 2^{52}\). Values fit in unsigned 64 bits and products for fraction comparisons fit in 128 bits. At width 32 a power-of-two shift supplies the rescaling; transitions are generated without storing a four-entry row for each cell. The two parities occupy disjoint locations. The old odd entry survives until its own ratio is tested.

\begin{center}
\small
\begin{tabular}{rrrr}
\toprule
Width & Stored cells & Round & Lower witness \(L(v)\), rounded down\\
\midrule
26 & 27\,088\,869 & 250 & 11.449834\\
28 & 54\,177\,740 & 300 & 11.541973\\
30 & 209\,295\,260 & 320 & 11.618175\\
32 & 810\,375\,650 & 300 & 11.680705\\
32 & 810\,375\,650 & 300, shifted & \result\\
\bottomrule
\end{tabular}
\end{center}
\noindent The first row retains all physical splits and tests both parities; its minimum also bounds the odd class. At width 32 the ordinary minimum exceeds \(11.68\), so \(w_a=2Av-23v\) and \(w_b=5Av-58v\) are strictly positive. Applying \(A\) once more gives their extrema exactly. The respective lower ratios are \(3074044368692/263088297614\) and \(18325770610324/1568251052756\). The second is the fraction in Theorem~\ref{main}. All \(386\,058\,577\) odd cells reach the root in exactly 30 single moves. This completes the proof. The smaller of the two upper ratios, \(15264493432827/1305926315603\), bounds only the width-32 growth constant.

\newpage
\section{A second certificate from one root}
Independently generate all physical matchings by the partner of the first end, followed by the interior and exterior matchings, and keep both banks distinct. At width 26 let \(A\) denote the two-step transfer. Write \(R_m=(A^m)_{oo}\). Concatenation, as in Lemma~\ref{strip}, gives \(\mu\ge R_m^{1/m}\) for every \(m\ge1\). No other entry of the final vector enters this inequality.

Set \(S=2^{45}\), \(x_0=S e_o\), choose integer shifts \(b_i\), and set \(x_i=\lfloor A x_{i-1}/2^{b_i}\rfloor\) entrywise. Then \(S A^m e_o\ge 2^{\sum b_i} x_m\). At \(m=3100\) the sum of the shifts is \(10894\) and \(x_m(o)=11\,568\,213\,483\,477\). The exact comparison
\[
 2^{10894}\cdot 11568213483477\cdot 10^{18600}
       > 2^{45}\cdot 11421465^{3100}
\]
proves \(\mu>11.421465\). This run used about 624 MiB of peak resident memory and 135.18 seconds. A second normalization, rounding down after multiplication by \(5/57\), first exceeds \(S\) at the root after 1859 slots and so gives \(\mu>11.4\) directly.

\paragraph{Replication.}
The accompanying archive contains the sources, integer logs, SHA-256 digests, and build and verification instructions. A direct enumeration of pairs of full matchings checks 42 combinations of width and order without using frontier transitions. A second coordinate scan decodes every matching and checks all \(209\,295\,260\) inserted pairs and \(552\,320\,023\) permitted splices at width 32; it also checks the \(209\,295\,261\) blocked splices where the neighboring ends were already paired. Another program reads the saved vector and computes both extrema directly with its own indexing and transfer code. For the two shifted vectors, \(A^2v\le 2^{60}{\bf1}\) bounds the unsigned 64-bit differences, and their cross products fit in 128 bits. Constructing the width-32 witness took 15 minutes 49 seconds with a peak resident set of 8\,948\,480 KiB on an Intel i9-13980HX. At width 26 the fully labelled graph has 27\,088\,869 states reachable from and able to return to the root, as checked by separate forward and reverse searches. At widths 26 through 32 the first twelve unrestricted root counts agree through \(M_{12}=678390116\). The folder for the independent return calculation includes both normalizations. The bounds at distinct widths refer to different subsets of meanders; the largest minimum gives a one-sided assertion about \(\mu\). Computing an individual unknown \(M_n\), or proving a limit for \(M_{n+1}/M_n\), requires additional information.

\newpage
\section{A prescribed relative error at every order}
The return calculation also supplies a finite denominator. In the width-32 comparison take the witness at round 180, with \(L(v)=2517567954515/215853621607>L=583/50\). Its coordinates are integers between \(1\) and \(2^{52}\), and the two-step adjacency \(A\) satisfies \(A^{15}e_o\ge{\bf1}\). Applying the inequality of Section 3 gives, for every \(n\ge16\),
\[
                 M_n\ge 2^{-52} L^{\,n-16}.
\]
This elementary prefactor suffices to turn a height estimate into an algorithm with a user-specified accuracy.

\begin{theorem}\label{accuracy}
For any \(n\ge1\) and rational \(0<\epsilon<1\), one can deterministically compute an integer \(H\) satisfying
\[
      (1-\epsilon)M_n\le H\le M_n
\]
in \(O^*(2^{4n/7}\epsilon^{-3/2})\) bit operations and bits of memory. Polynomial factors in \(n\) and the input bit length are suppressed.
\end{theorem}
\begin{proof}
Within a factor, match brackets separately on each bank, leaving unmatched ends free. A cycle in the union of the internal pairs persists in every extension: each of its sites already has both incident edges. This includes a two-edge cycle when both banks pair the same sites. Such a factor cannot occupy a proper interval of a connected meander. Retain the 1\,207 cycle-containing factors of length at most ten with no proper cycle-containing subfactor. After each letter, reject a completed forbidden suffix; otherwise keep the longest suffix that is a proper prefix of a forbidden factor (possibly empty). This gives 2\,899 states. Weight each simultaneous opening by \(Q=159/100\), each simultaneous closing by \(1/Q\), and either shift by one. For its accepted-letter matrix \(D\), put \(s=(15900D)^{48}{\bf1}\). The integer checks give \(Ds\le r s\), \(r=3898274927/10^9\), and \(s_{\rm empty}^2<2(\min s)^2\). Hence the total weight of accepted words of length \(t\) is at most \((s_{\rm empty}/\min s)r^t\). The companion manuscript \cite{halves} gives the factor test and finite counting details.

Consider a cut of height \(h\). Read each half from its outside end, taking the reversed bracket directions for the right half. Each half carries weight \(Q^{h/2}\). Dropping the requirement that their boundary pairings glue and summing over the \(2n-1\) cuts bounds the number of meanders omitted by width \(W\) by \(2(2n-1)r^{2n}/Q^W\). For \(n\ge16\) take the smallest even \(W\) satisfying
\[
 Q^W\ \ge\ \frac{2^{53}(2n-1)r^{2n}}{\epsilon L^{\,n-16}},
\]
or take \(W=2n\) if this is smaller. For \(n<16\) use full width. The integer \(H=C_n(W)\) then has the claimed relative accuracy.

The comparisons \((r^2/L)^7<Q^4\) and \(Q^3>4\) are rational. The first gives \(W\le 4n/7+\log_Q(n/\epsilon)+O(1)\). The strip has \(O^*(2^W)\) states with polynomial bit cost per transition and \(O(n)\)-bit counts. As \(\log_Q2<3/2\), the second comparison yields the stated bound. All width choices may be made by exact integer comparisons.
\end{proof}

\newpage
\section{Finite constants and checks on the approximation}
A single test of the positive-vector inequality also gives a lower bound at every finite order. For a witness \(w>0\) with \(Aw\ge\lambda w\), compute \(y\) and an integer \(E\) with \(2^E y\le A^b e_o\). If \(g=\min_i y_i/w_i>0\), positivity gives
\[
  M_n\ \ge\ 2^E\,g\,w_o\,\lambda^{n-1-b}\qquad(n\ge b+1).
\]
For \(n\le b\) one can instead use the root entries already recorded while constructing \(y\). This avoids fitting a prefactor to the known exact meander numbers. The supplement carries out both comparisons at width 32. For \(w=5Av-58v\) it gives \(w_o=107564340496998\), and at \(b=40\) the recorded exponent and minimum ratio are \(E=76\) and \(g=308588629582/89512920981\). Comparing the forty earlier root entries as well gives the uniform bounds, for all \(n\ge1\),
\[
 M_n>4.71986\cdot 10^{-7}\,(\wnum/\wden)^n,\qquad
 M_n\ge 2^{-21}(11.68)^n.
\]

For the height calculation, counting restricted Dyck pairs by their two nonnegative heights and the avoidance state can give a sharper certificate at a particular order. With factors through length eight let \(B_n(W)\) count all such pairs staying within width \(W\). Every omitted connected meander belongs to the difference, so
\[
 0\le M_n-C_n(W)\le B_n(2n)-B_n(W)\quad(n>4).
\]
A fixed automaton takes polynomial time in \(n\) for this difference. The companion computations bound the relative loss at \(W=32\) by \(0.015420\%\) for order 30, \(0.067896\%\) for order 32 and \(3.036153\%\) for order 39 (rounded up). The first two differences can be tested against their known exact counts. At order 39 the earlier independently certified interval in the companion data is narrower.

For the accuracy rule as printed in Theorem~\ref{accuracy}, the checker returns \(W=762\) at \(n=1000,\epsilon=1/100\), and \(W=782\) at the same order with \(\epsilon=10^{-6}\). Using the uniform denominator above reduces these widths to 628 and 648. These are sufficient width calculations; enumeration at order 1000 is unnecessary to check them. A reference implementation evaluates the restricted count itself for tractable inputs, for example \(H=1828\) at order six.

\newpage
\section*{Data and code availability}
The exact integer certificates, verifier sources, companion manuscript and reproduction instructions are archived at
\href{https://doi.org/10.5281/zenodo.24448957}{doi:10.5281/zenodo.24448957}.
The \href{https://github.com/Kimargin8/meander-growth-bound}{source repository on GitHub} also provides the verification programs.

\section*{Declaration of generative AI use}
The author used OpenAI ChatGPT and Codex to explore mathematical arguments, develop and test the verification programs, and draft and revise the text. The integer witnesses and the programs used to check them are included in the accompanying archive.

\begin{thebibliography}{9}\small
\bibitem{halves} I.~Argunov, \emph{Shortest and near-shortest halves of closed plane meanders}, companion manuscript (2026), sections on simultaneous cuts and approximation.
\bibitem{ap} M.~H. Albert and M.~S. Paterson, \emph{Bounds for the growth rate of meander numbers}, J. Combin. Theory Ser. A \textbf{112} (2005), 250--262.
\href{https://doi.org/10.1016/j.jcta.2005.02.006}{Journal article}.
\bibitem{bacher} R. Bacher, \emph{Meander algebras}, Institut Fourier.
\href{https://www-fourier.univ-grenoble-alpes.fr/sites/default/files/ref_478.pdf}{Author's manuscript}.
\bibitem{bfw} J.~Borga, V.~F\'eray and P.~Winkler, \emph{Permutons} (2026).
\href{https://arxiv.org/abs/2611.12104}{arXiv:2611.12104}.
\bibitem{caroppo} S.~Caroppo, G.~Da Lozzo, G.~Di Battista and J.~Vihrovs, \emph{Faster planar graph algorithms for connectivity problems via meanders} (2026).
\href{https://arxiv.org/abs/2610.11944}{arXiv:2610.11944}.
\bibitem{jensen} I.~Jensen, \emph{A transfer matrix approach to the enumeration of plane meanders}, J. Phys. A (2000).
\href{https://doi.org/10.1088/0305-4470/33/34/301}{Journal article}.
\end{thebibliography}
\end{document}
